% =========================================================================
% Capítulo 3: Cinética de sistemas de partículas
% Tema 01: Sistemas de partículas
% Problema 02: Péndulo de mancuerna en forma de T
% Origen: Ayudantía 29, 2026
% =========================================================================
\subsubsection{Péndulo de mancuerna en forma de T}

Una mancuerna de masa $2m$ y largo $2a$, compuesta por dos partículas de masa $m$ en sus extremos, está soldada a una barra rígida sin masa de largo $l$ en su punto medio $P$, formando una estructura en forma de letra $T$. El sistema es libre de rotar en el plano vertical alrededor de una articulación fija en el punto $O$, como se muestra en la \textbf{Figura \ref{fig:mancuerna_T}}. Sabiendo que la gravedad actúa hacia abajo y que el ángulo $\theta$ se mide desde la vertical descendente hacia la barra de longitud $l$, realice lo siguiente:

\begin{enumerate}
    \item Defina los marcos de referencia y sistemas de coordenadas adecuados.
    \item Analice la cinemática del cuerpo rígido, determinando el momento de inercia del sistema respecto al punto de rotación $O$.
    \item Analice la cinética del sistema para obtener la ecuación diferencial de movimiento en términos de $\theta$ utilizando balance de momentum angular y el teorema de trabajo y energía.
\end{enumerate}

\begin{figure}[!ht]
    \centering
    \begin{tikzpicture}[>=Latex, scale=1.3]
        % --- Parámetros ---
        \def\ang{30}      % Ángulo theta
        \def\Lbar{3.0}    % Longitud l
        \def\Aarm{1.2}    % Mitad del largo de la mancuerna (a)

        % --- Soporte Fijo O ---
        \fill[pattern=north east lines] (-0.5, 0.1) rectangle (0.5, 0.4);
        \draw[thick] (-0.5, 0.1) -- (0.5, 0.1);
        \draw[thick, fill=gray!20] (-0.2,0.1) arc (0:180:-0.2) -- cycle;
        
        \coordinate (O) at (0,0);
        \filldraw (O) circle (2pt) node[left, xshift=-5pt] {$O$};

        % --- Barra y Mancuerna ---
        \begin{scope}[rotate=-\ang]
            % Barra principal l
            \draw[line width=4pt] (0,0) -- (0, -\Lbar);
            \draw[line width=4pt] (-\Aarm, -\Lbar) -- (\Aarm, -\Lbar);
            \draw[line width=3pt, gray!40, line cap=round] (0,0) -- (0, -\Lbar);
            \filldraw (O) circle (0.6pt);
            %\node[left] at (0, -\Lbar/2) {$l$};
            
            \coordinate (P) at (0, -\Lbar);
            \node[below left] at (P) {$P$};

            % Travesaño de la mancuerna (2a)
            \draw[line width=3pt, gray!40] (-\Aarm, -\Lbar) -- (\Aarm, -\Lbar);
            \draw[thick] (0,0) -- (0, -\Lbar);
            \draw[thick] (-\Aarm, -\Lbar) -- (\Aarm, -\Lbar);
            \filldraw (P) circle (0.6pt);
            
            % Masas m
            \shadedraw[ball color=gray!30] (-\Aarm, -\Lbar) circle (0.2) node[font=\footnotesize] {$m$};
            \shadedraw[ball color=gray!30] (\Aarm, -\Lbar) circle (0.2) node[font=\footnotesize] {$m$};
            
            % Cota 2a
            \draw[|<->|] (-\Aarm, -\Lbar-0.5) -- (\Aarm, -\Lbar-0.5) node[midway, below] {$2a$};
            
            % Cota l (indicador lateral)
            \draw[|<->|] (-0.8, 0) -- (-0.8, -\Lbar) node[midway, above left] {$l$};
        \end{scope}

        % --- Referencias ---
        % Línea vertical de referencia
        \draw[dashed, gray] (O) -- (0, -2.5);
        
        % Ángulo theta
        \draw[->] (0, -1.0) arc (-90:{-90-\ang}:1.0);
        \node at (-0.3, -1.3) {$\theta$};

        % Vector Gravedad
        \draw[->, thick] (1.5, -0.5) -- (1.5, -1.5) node[midway, right] {$\mathbf{g}$};

    \end{tikzpicture}
    \caption{Esquema de una mancuerna rígida en rotación suspendida de un punto fijo $O$.}
    \label{fig:mancuerna_T}
\end{figure}

% --- SOLUCIÓN ---------------------------------------------------------------
\begin{solucion}{Péndulo de mancuerna en forma de T (por momento angular)}
\textbf{Marcos de referencia y Sistemas de Coordenadas}
\begin{itemize}
    \item $\mathcal{F}$: inercial $\longrightarrow\{\vec{e}_x;\vec{e}_y;\vec{e}_z\}$
    \item $\mathcal{A}$: ligado a la barra, con rotación en $\vec{e}_z$ $\longrightarrow\{\vec{e}_r;\vec{e}_\theta;\vec{e}_z\}$
\end{itemize}

\begin{center}
\begin{tikzpicture}[>=Latex, scale=1.2]
    % --- PARÁMETROS ---
    \def\ang{30}   % ángulo theta [°]
    \def\len{1.6}  % largo de los vectores [u]

    % --- BASE FIJA (e_x hacia abajo, e_y hacia la izquierda, e_z entrando) ---
    \draw[->, thick, azul] (0,0) -- (0,-\len) node[below] {$\vec{e}_x$};
    \draw[->, thick, azul] (0,0) -- (-\len,0) node[left] {$\vec{e}_y$};
    \draw[gray] (0,0) circle (0.12);
    \draw[gray] (-0.085,-0.085) -- (0.085,0.085) (-0.085,0.085) -- (0.085,-0.085);
    \node[gray, above right] at (0.05,0.12) {$\vec{e}_z$};

    % --- BASE MÓVIL ---
    \draw[->, thick, myred] (0,0) -- ({-90-\ang}:\len) node[below left] {$\vec{e}_r$};
    \draw[->, thick, myred] (0,0) -- ({180-\ang}:\len) node[above left] {$\vec{e}_\theta$};

    % --- ÁNGULO ---
    \draw[->] (0,-0.8) arc (-90:{-90-\ang}:0.8);
    \node at ({-90-\ang/2}:1.05) {$\theta$};
    \node[below right] at (0.05,-0.05) {$O$};
\end{tikzpicture}
\end{center}

Relaciones cinemáticas en coordenadas polares (en el plano $x\!-\!y$):
$$\dot{\vec{e}}_r=\dot\theta\,\vec{e}_\theta,\qquad \dot{\vec{e}}_\theta=-\dot\theta\,\vec{e}_r,\qquad \vec{e}_x=\cos\theta\,\vec{e}_r-\sin\theta\,\vec{e}_\theta .$$
Geometría de las masas (dos partículas de masa $m$):
$$\vec{r}_1= \ell\,\vec{e}_r - a\,\vec{e}_\theta, \qquad \vec{r}_2= \ell\,\vec{e}_r + a\,\vec{e}_\theta .$$

\subsubsection*{Análisis cinemático}
Velocidades (derivando en $\mathcal{F}$ y usando $\dot{\vec{e}}_r,\dot{\vec{e}}_\theta$):
$$\begin{aligned}
\vec{v}_1 &= \frac{d}{dt}\big(\ell\,\vec{e}_r - a\,\vec{e}_\theta\big) = \ell\dot\theta\,\vec{e}_\theta + a\dot\theta\,\vec{e}_r,\\[2mm]
\vec{v}_2 &= \frac{d}{dt}\big(\ell\,\vec{e}_r + a\,\vec{e}_\theta\big) = \ell\dot\theta\,\vec{e}_\theta - a\dot\theta\,\vec{e}_r.
\end{aligned}$$

\subsubsection*{Análisis cinético}
Momento angular respecto de $O$ en el marco $\mathcal{F}$:
$$\begin{aligned}
\marc{\vec{H}_O}{F}{}
&= \sum_{i=1}^{2} \vec{r}_i \times m\,\vec{v}_i \\[1mm]
&= (\ell\,\vec{e}_r - a\,\vec{e}_\theta)\times m(\ell\dot\theta\,\vec{e}_\theta + a\dot\theta\,\vec{e}_r)
 \;+\; (\ell\,\vec{e}_r + a\,\vec{e}_\theta)\times m(\ell\dot\theta\,\vec{e}_\theta - a\dot\theta\,\vec{e}_r)\\[1mm]
&= m\dot\theta\Big[\ell^2\,\vec{e}_r\!\times\!\vec{e}_\theta + a^2\,\vec{e}_\theta\!\times\!\vec{e}_r\Big]
 + m\dot\theta\Big[\ell^2\,\vec{e}_r\!\times\!\vec{e}_\theta + a^2\,\vec{e}_\theta\!\times\!\vec{e}_r\Big] \\[1mm]
&= 2m\dot\theta\,(\ell^2+a^2)\,\vec{e}_z .
\end{aligned}$$
De aquí se lee directamente el momento de inercia respecto de $O$, ya que $\marc{\vec{H}_O}{F}{}=I_O\dot\theta\,\vec{e}_z$:
$$I_O = 2m(\ell^2+a^2)$$
Como $O$ es fijo en el inercial,
$$\der{\marc{\vec{H}_O}{F}{}}{F}{} = 2m\ddot\theta\,(\ell^2+a^2)\,\vec{e}_z.$$

Diagrama de cuerpos libres: las únicas fuerzas externas relevantes son los pesos $m\vec{g}$ de cada partícula. Sus momentos respecto de $O$:
$$\marc{\vec{M}_O}{F}{} = \vec{r}_1\times m\vec{g} + \vec{r}_2\times m\vec{g}.$$
Tomando $\vec{g} = -g\,\vec{e}_x$ y usando $\vec{e}_x=\cos\theta\,\vec{e}_r-\sin\theta\,\vec{e}_\theta$:
$$\begin{aligned}
\vec{r}_1\times m\vec{g}
&= (\ell\,\vec{e}_r - a\,\vec{e}_\theta)\times\big[-mg(\cos\theta\,\vec{e}_r-\sin\theta\,\vec{e}_\theta)\big]
= mg(-\ell\sin\theta + a\cos\theta)\,\vec{e}_z,\\[1mm]
\vec{r}_2\times m\vec{g}
&= (\ell\,\vec{e}_r + a\,\vec{e}_\theta)\times\big[-mg(\cos\theta\,\vec{e}_r-\sin\theta\,\vec{e}_\theta)\big]
= mg(-\ell\sin\theta - a\cos\theta)\,\vec{e}_z.
\end{aligned}$$
Por lo tanto,
$$\marc{\vec{M}_O}{F}{}
= \big[mg(-\ell\sin\theta + a\cos\theta) + mg(-\ell\sin\theta - a\cos\theta)\big]\vec{e}_z
= -2mg\ell\sin\theta\,\vec{e}_z .$$

\subsubsection*{Ecuación de movimiento}
Balance de momento angular en $O$:
$$\der{\marc{\vec{H}_O}{F}{}}{F}{} = \marc{\vec{M}_O}{F}{}
\quad\Longrightarrow\quad
2m\ddot\theta(\ell^2+a^2)\,\vec{e}_z = -2mg\ell\sin\theta\,\vec{e}_z.$$
Cancelando factores comunes y proyectando en $\vec{e}_z$:
\respuesta{$\displaystyle \ddot\theta + \frac{g\ell}{\ell^2+a^2}\,\sin\theta = 0$}

\patotruco{Cuando todo el cuerpo gira alrededor de un punto fijo $O$, el balance de momento angular se reduce a $I_O\ddot\theta=M_O$. Aquí $I_O=\sum m_i d_i^2$ con $d_i^2=\ell^2+a^2$ para ambas masas: nos ahorramos todos los productos cruz.}
\end{solucion}

\begin{solucion}{Péndulo de mancuerna en forma de T (por trabajo y energía)}
\textbf{Balance por teorema de trabajo y energía}\\
Con los mismos marcos y las mismas velocidades $\vec{v}_1,\vec{v}_2$ de la solución anterior, la energía cinética total en $\mathcal{F}$ es
$$\marc{T}{F}{}=\sum_{i=1}^{2}\frac{1}{2}\,m\,\vec{v}_i\!\cdot\!\vec{v}_i .$$
Teorema de trabajo y energía en $\mathcal{F}$: la variación de la energía cinética es la potencia de las fuerzas internas y externas,
$$\frac{d}{dt}\big(\marc{T}{F}{}\big)=\sum_{i=1}^{2}\vec{R}_i\!\cdot\!\vec{v}_i .$$

\textbf{Cálculo de la energía cinética}\\
Con $\vec{v}_1=\ell\dot\theta\,\vec{e}_\theta + a\dot\theta\,\vec{e}_r$ y $\vec{v}_2=\ell\dot\theta\,\vec{e}_\theta - a\dot\theta\,\vec{e}_r$,
$$\begin{aligned}
\marc{T}{F}{}
&= \tfrac12 m\,(\ell\dot\theta\,\vec{e}_\theta + a\dot\theta\,\vec{e}_r)\!\cdot\!(\ell\dot\theta\,\vec{e}_\theta + a\dot\theta\,\vec{e}_r)
   + \tfrac12 m\,(\ell\dot\theta\,\vec{e}_\theta - a\dot\theta\,\vec{e}_r)\!\cdot\!(\ell\dot\theta\,\vec{e}_\theta - a\dot\theta\,\vec{e}_r)\\[1mm]
&= \tfrac12 m\,\dot\theta^2\,[\ell^2 + a^2] + \tfrac12 m\,\dot\theta^2\,[\ell^2 + a^2]
= m\dot\theta^2\,[\ell^2+a^2].
\end{aligned}$$
Por lo tanto,
$$\frac{d}{dt}\,\marc{T}{F}{} = 2\,\dot\theta\,\ddot\theta\,m\,[\ell^2+a^2].$$

\textbf{Cálculo de la potencia de las partículas}\\
Sean $\vec{r}_1=\ell\,\vec{e}_r-a\,\vec{e}_\theta$ y $\vec{r}_2=\ell\,\vec{e}_r+a\,\vec{e}_\theta$. La única fuerza externa que trabaja sobre cada partícula es el peso, con $\vec{g}=-g\,\vec{e}_x$ y $\vec{e}_x=\cos\theta\,\vec{e}_r-\sin\theta\,\vec{e}_\theta$. Así,
$$\sum_{i=1}^{2}\vec{R}_i\!\cdot\!\vec{v}_i
= m g(\cos\theta\,\vec{e}_r-\sin\theta\,\vec{e}_\theta)\!\cdot\!(\ell\dot\theta\,\vec{e}_\theta + a\dot\theta\,\vec{e}_r)
 + m g(\cos\theta\,\vec{e}_r-\sin\theta\,\vec{e}_\theta)\!\cdot\!(\ell\dot\theta\,\vec{e}_\theta - a\dot\theta\,\vec{e}_r).$$
Evaluando los productos escalares:
$$\sum_{i=1}^{2}\vec{R}_i\!\cdot\!\vec{v}_i
= mg\,[a\dot\theta\cos\theta-\ell\dot\theta\sin\theta]
 + mg\,[-a\dot\theta\cos\theta-\ell\dot\theta\sin\theta]
= 2mg\,\ell\,\dot\theta\,\sin\theta .$$

\textbf{Igualando}\\
Aplicando el teorema de trabajo y energía:
$$2\,\dot\theta\,\ddot\theta\,m\,[\ell^2+a^2] \;=\; -\,2\,m g\,\ell\,\dot\theta\,\sin\theta.$$
Si $\dot\theta\neq 0$, se simplifica y queda la ecuación de movimiento:
\respuesta{$\displaystyle \ddot\theta+\frac{g\,\ell}{\ell^2+a^2}\,\sin\theta=0$}

\patoojo{Este método solo da \emph{una} ecuación escalar, que es justo lo que necesitamos porque el sistema tiene un único grado de libertad. Pero ojo con el paso final: dividimos por $\dot\theta$, así que la ecuación resultante vale también cuando $\dot\theta=0$ solo porque coincide con la del balance de momento angular.}
\end{solucion}
